\textbf{Phần 1. Trắc nghiệm lựa chọn}
\setcounterpageref{ex}{0}

\begin{ex} %[5H9N53]
Muốn tính thể tích hình lập phương có cạnh $a$, ta dùng công thức
\choice
{$V = a \times 4$}
{\True $V = a \times a \times a$}
{$V = a \times a \times 6$}
{$V = a \times a \times 4$}
\loigiai{
$V = a \times a \times a$.\\
\textbf{Đáp án của câu hỏi là B}
}
\end{ex}

\begin{ex} %[5H9N53]
Một hình lập phương có cạnh $4\text{ cm}$. Thể tích của hình lập phương đó là
\choice
{$16\text{ cm}^3$}
{$96\text{ cm}^3$}
{\True $64\text{ cm}^3$}
{$64\text{ cm}^2$}
\loigiai{
$V = 4 \times 4 \times 4 = 64\text{ (cm}^3\text{)}$.\\
\textbf{Đáp án của câu hỏi là C}
}
\end{ex}

\begin{ex} %[5H9H53]
Một khối gỗ hình lập phương có cạnh $0{,}5\text{ m}$. Thể tích khối gỗ đó là
\choice
{$1{,}25\text{ m}^3$}
{\True $0{,}125\text{ m}^3$}
{$125\text{ dm}^3$}
{$1{,}5\text{ m}^3$}
\loigiai{
$V = 0{,}5 \times 0{,}5 \times 0{,}5 = 0{,}125\text{ (m}^3\text{)}$.\\
\textbf{Đáp án của câu hỏi là B}
}
\end{ex}

\begin{ex} %[5H9H53]
Một hình lập phương có diện tích một mặt là $25\text{ cm}^2$. Thể tích của hình lập phương đó là
\choice
{$100\text{ cm}^3$}
{$150\text{ cm}^3$}
{$625\text{ cm}^3$}
{\True $125\text{ cm}^3$}
\loigiai{
Vì $S_{1mặt} = 25\text{ cm}^2$ nên cạnh $a = 5\text{ cm}$.\\
$V = 5 \times 5 \times 5 = 125\text{ (cm}^3\text{)}$.\\
\textbf{Đáp án của câu hỏi là D}
}
\end{ex}

\begin{ex} %[5H9H53]
Một hình lập phương có diện tích toàn phần là $54\text{ dm}^2$. Thể tích của hình lập phương đó là
\choice
{\True $27\text{ dm}^3$}
{$18\text{ dm}^3$}
{$36\text{ dm}^3$}
{$81\text{ dm}^3$}
\loigiai{
Diện tích một mặt: $54 : 6 = 9\text{ (dm}^2\text{)}$. Cạnh là $3\text{ dm}$.\\
$V = 3 \times 3 \times 3 = 27\text{ (dm}^3\text{)}$.\\
\textbf{Đáp án của câu hỏi là A}
}
\end{ex}

\begin{ex} %[5H9H53]
Một thùng chứa dạng hình lập phương có cạnh $1\text{ m}$. Thùng đó chứa được tối đa bao nhiêu lít nước?
\choice
{$100\text{ l}$}
{$10\,000\text{ l}$}
{\True $1\,000\text{ l}$}
{$100\,000\text{ l}$}
\loigiai{
$V = 1 \times 1 \times 1 = 1\text{ m}^3 = 1\,000\text{ dm}^3 = 1\,000\text{ l}$.\\
\textbf{Đáp án của câu hỏi là C}
}
\end{ex}

\begin{ex} %[5H9V53]
Nếu tăng cạnh của một hình lập phương lên gấp $2$ lần thì thể tích của nó tăng lên gấp mấy lần?
\choice
{$2$ lần}
{\True $8$ lần}
{$4$ lần}
{$6$ lần}
\loigiai{
Thể tích mới: $V' = (2a) \times (2a) \times (2a) = 8 \times a^3 = 8V$.\\
\textbf{Đáp án của câu hỏi là B}
}
\end{ex}

\begin{ex} %[5H9V53]
Một khối sắt hình lập phương có cạnh $0{,}2\text{ m}$. Biết $1\text{ dm}^3$ sắt nặng $7{,}8\text{ kg}$. Khối sắt đó nặng bao nhiêu ki-lô-gam?
\choice
{\True $62{,}4\text{ kg}$}
{$624\text{ kg}$}
{$6{,}24\text{ kg}$}
{$31{,}2\text{ kg}$}
\loigiai{
$V = 0{,}2 \times 0{,}2 \times 0{,}2 = 0{,}008\text{ m}^3 = 8\text{ dm}^3$.\\
Khối lượng: $8 \times 7{,}8 = 62{,}4\text{ (kg)}$.\\
\textbf{Đáp án của câu hỏi là A}
}
\end{ex}

\begin{ex} %[5H9V53]
Người ta xếp các khối lập phương nhỏ cạnh $1\text{ cm}$ thành một hình lập phương lớn có cạnh $6\text{ cm}$. Số khối lập phương nhỏ cần dùng là
\choice
{$36$ khối}
{$144$ khối}
{$108$ khối}
{\True $216$ khối}
\loigiai{
Số khối: $6 \times 6 \times 6 = 216$ khối.\\
\textbf{Đáp án của câu hỏi là D}
}
\end{ex}

\begin{ex} %[5H9V53]
Hai hình lập phương có tỉ số cạnh là $\dfrac{1}{3}$. Tỉ số thể tích của hình bé so với hình lớn là
\choice
{\True $\dfrac{1}{27}$}
{$\dfrac{1}{3}$}
{$\dfrac{1}{9}$}
{$\dfrac{1}{6}$}
\loigiai{
Tỉ số thể tích: $\left(\dfrac{1}{3}\right)^3 = \dfrac{1}{27}$.\\
\textbf{Đáp án của câu hỏi là A}
}
\end{ex}

\begin{ex} %[5H9V53]
Một hình lập phương có diện tích xung quanh là $64\text{ cm}^2$. Thể tích của nó là
\choice
{$32\text{ cm}^3$}
{$96\text{ cm}^3$}
{$128\text{ cm}^3$}
{\True $64\text{ cm}^3$}
\loigiai{
Diện tích một mặt: $64 : 4 = 16\text{ (cm}^2\text{)}$. Cạnh là $4\text{ cm}$.\\
$V = 4 \times 4 \times 4 = 64\text{ (cm}^3\text{)}$.\\
\textbf{Đáp án của câu hỏi là D}
}
\end{ex}

\begin{ex} %[5H9C53]
Một khối gỗ hình lập phương có cạnh $9\text{ cm}$. Người ta khoét một lỗ rỗng cũng hình lập phương có cạnh $3\text{ cm}$ ở chính giữa khối gỗ. Thể tích phần gỗ còn lại là
\choice
{$720\text{ cm}^3$}
{$726\text{ cm}^3$}
{\True $702\text{ cm}^3$}
{$648\text{ cm}^3$}
\loigiai{
Thể tích ban đầu: $9 \times 9 \times 9 = 729\text{ (cm}^3\text{)}$.\\
Thể tích phần khoét: $3 \times 3 \times 3 = 27\text{ (cm}^3\text{)}$.\\
Thể tích còn lại: $729 - 27 = 702\text{ (cm}^3\text{)}$.\\
\textbf{Đáp án của câu hỏi là C}
}
\end{ex}

\textbf{Phần 2. Trắc nghiệm đúng sai}
\setcounterpageref{ex}{0}

\begin{ex} %[5H9H53]
Cho hình lập phương có cạnh $6\text{ dm}$. Xét tính đúng sai:
\choiceTF
{\True Diện tích toàn phần của hình lập phương là $216\text{ dm}^2$.}
{\True Thể tích của hình lập phương là $216\text{ dm}^3$.}
{Số đo diện tích toàn phần (theo $\text{dm}^2$) và thể tích (theo $\text{dm}^3$) của hình này luôn khác nhau với mọi hình lập phương.}
{\True Với hình lập phương cạnh $6\text{ dm}$, số đo diện tích toàn phần và số đo thể tích có cùng giá trị số là $216$.}
\loigiai{
\begin{itemize}
    \item $S_{tp} = 6 \times 6 \times 6 = 216\text{ dm}^2$. Đúng.
    \item $V = 6 \times 6 \times 6 = 216\text{ dm}^3$. Đúng.
    \item Cạnh bằng 6 thì giá trị số trùng nhau nên khẳng định luôn khác nhau là Sai.
    \item Đúng, cùng giá trị 216.
\end{itemize}
\begin{itemize}
    \item a) \textbf{Đúng}.
    \item b) \textbf{Đúng}.
    \item c) \textbf{Sai}.
    \item d) \textbf{Đúng}.
\end{itemize}
}
\end{ex}

\begin{ex} %[5H9V53]
Một bể nước hình lập phương có cạnh trong lòng bể là $1{,}2\text{ m}$. Bể đang chứa $864\text{ l}$ nước. Xét tính đúng sai:
\choiceTF
{\True Dung tích toàn bộ của bể là $1\,728\text{ l}$.}
{\True Lượng nước hiện có trong bể chiếm $50\%$ thể tích bể.}
{Mực nước trong bể hiện cao $0{,}8\text{ m}$.}
{\True Mực nước trong bể hiện cao $0{,}6\text{ m}$.}
\loigiai{
\begin{itemize}
    \item Dung tích: $1{,}2 \times 1{,}2 \times 1{,}2 = 1{,}728\text{ m}^3 = 1728\text{ l}$. Đúng.
    \item Tỉ lệ nước: $864 : 1728 = 50\%$. Đúng.
    \item Chiều cao nước: $1{,}2 \times 50\% = 0{,}6\text{ m} \neq 0{,}8\text{ m}$. Ý 3 Sai, Ý 4 Đúng.
\end{itemize}
\begin{itemize}
    \item a) \textbf{Đúng}.
    \item b) \textbf{Đúng}.
    \item c) \textbf{Sai}.
    \item d) \textbf{Đúng}.
\end{itemize}
}
\end{ex}

\textbf{Phần 3. Trả lời ngắn}
\setcounterpageref{ex}{0}

\begin{ex} %[5H9H53]
Một hình lập phương có cạnh $8\text{ cm}$. Tính thể tích theo xăng-ti-mét khối.
\shortans{512}
\loigiai{
\textbf{Đáp số:} $512$\\
$8 \times 8 \times 8 = 512\text{ cm}^3$.
}
\end{ex}

\begin{ex} %[5H9H53]
Một hình lập phương có diện tích một mặt là $49\text{ cm}^2$. Tính thể tích hình lập phương đó theo xăng-ti-mét khối.
\shortans{343}
\loigiai{
\textbf{Đáp số:} $343$\\
Cạnh là $7\text{ cm} \Rightarrow V = 7 \times 7 \times 7 = 343\text{ cm}^3$.
}
\end{ex}

\begin{ex} %[5H9V53]
Một thùng chứa hình lập phương có cạnh $1{,}5\text{ m}$. Tính số lít nước tối đa thùng chứa được.
\shortans{3375}
\loigiai{
\textbf{Đáp số:} $3375$\\
$1{,}5 \times 1{,}5 \times 1{,}5 = 3{,}375\text{ m}^3 = 3\,375\text{ l}$.
}
\end{ex}

\begin{ex} %[5H9V53]
Một khối gỗ hình lập phương có thể tích $1\,000\text{ cm}^3$. Tính diện tích toàn phần của khối gỗ đó theo xăng-ti-mét vuông.
\shortans{600}
\loigiai{
\textbf{Đáp số:} $600$\\
Cạnh: $10\text{ cm} \Rightarrow S_{tp} = 10 \times 10 \times 6 = 600\text{ cm}^2$.
}
\end{ex}

\begin{ex} %[5H9V53]
Một khối nhôm hình lập phương cạnh $4\text{ dm}$. Biết $1\text{ dm}^3$ nhôm nặng $2{,}7\text{ kg}$. Tính khối lượng khối nhôm theo ki-lô-gam.
\shortans{172,8}
\loigiai{
\textbf{Đáp số:} $172{,}8$\\
Thể tích: $4 \times 4 \times 4 = 64\text{ dm}^3$.\\
Khối lượng: $64 \times 2{,}7 = 172{,}8\text{ kg}$.
}
\end{ex}

\begin{ex} %[5H9C53]
Có một khối lập phương bằng kim loại cạnh $12\text{ cm}$. Người ta đem nung chảy rồi đúc thành các khối lập phương nhỏ cạnh $3\text{ cm}$. Hỏi đúc được bao nhiêu khối lập phương nhỏ như vậy?
\shortans{64}
\loigiai{
\textbf{Đáp số:} $64$\\
$(12 \times 12 \times 12) : (3 \times 3 \times 3) = 1\,728 : 27 = 64$ khối.
}
\end{ex}
